\subsection*{Переход от квантовых фононов C46 к классическим}

Сохранены авторские свободные энергии 46 структур на девяти объёмах и
температурах от 0 до 2200~К с шагом 25~К. Полный комплект исходных
спектров для этих структур отсутствует в архиве. Для каждой устойчивой
пары «структура~--- объём» из полной энергии вычли статическую метку
и восстановили классическую колебательную энергию при \(T_0=1000\)~К.
Для одной гармонической моды частоты \(\omega\) при
\(z=\hbar\omega/(k_{\rm B}T)\) квантовая свободная энергия равна
\(k_{\rm B}T\ln[2\sinh(z/2)]\), классическая~--- \(k_{\rm B}T\ln z\).
Их разность начинается с \(k_{\rm B}T(z^2/24-z^4/2880)\).
После суммирования по модам и нормировки на атом поэтому получаем
\[
 \frac{F_{\rm q}^{\rm vib}(T)}{T/T_0}
 = F_{\rm cl}^{\rm vib}(T_0)
   -g k_{\rm B}T_0\ln(T/T_0)
   +\sum_{j=1}^{4} b_j(T_0/T)^{2j}
   +O\big((T_0/T)^{10}\big).
\]
Искомая энергия~--- свободный член. Вес \(g\) сохранённых мод на атом
определяли вместе с коэффициентами \(b_j\), чтобы учесть обработку
поступательных мод в авторской фононной сетке. При подгонке на
1000--2200~К и проверке с нижними границами 800, 1200 и 1400~К и
порядками 5, 4 и 3 максимальный разброс \(F_{\rm cl}^{\rm vib}\)
составил \(2{,}4\cdot10^{-12}\)~эВ/атом по всем 409 устойчивым парам.
Для 396 устойчивых спектров из независимой выборки ATAT ту же процедуру
проверили прямой суммой по частотам. Максимальная ошибка составила
\(7{,}0\cdot10^{-13}\)~эВ/атом. Восемь химических заполнений совпали
с авторскими. По всем девяти объёмам среднеквадратическое расхождение
между прямыми классическими энергиями ATAT и реконструкцией C46 для этих
структур составляет \(0{,}0136\)~мэВ/атом, при \(V_0\)~---
\(0{,}00104\)~мэВ/атом.

На каждом объёме классические метки подогнали в исходном базисе из
20 орбит с прежними весами, масками обучения и устойчивости.
Статические взаимодействия, отношение \(c/a=1{,}56374\) и обработка
объёмной зависимости сохранены. Поправка
\(\delta f=F_{\rm cl}^{\rm vib}-F_{\rm q}^{\rm vib}\) по всей выборке
лежит между \(-1{,}217\) и \(-0{,}491\)~мэВ/атом. При опорном объёме
её наилучшая линейная по составу часть по 46 структурам составляет
\(-0{,}7212-0{,}6661x_{\rm Al}\)~мэВ/атом, где \(x_{\rm Al}\)~---
атомная доля. Среднеквадратический остаток равен \(0{,}0466\), максимальный
по модулю~--- \(0{,}1251\)~мэВ/атом. При фиксированном объёме линейный
по составу член меняет наклон общей касательной и сохраняет составы.
В QHA все члены удержаны, поскольку их зависимость от объёма влияет на
минимум свободной энергии.

\begin{center}
\small
\begin{tabular}{@{}lrrrrrr@{}}
\toprule
 & \multicolumn{2}{c}{Квантовые, ат.\% Al}
 & \multicolumn{2}{c}{Классические, ат.\% Al}
 & \multicolumn{2}{c}{Сдвиг, п.п.}\\
\cmidrule(lr){2-3}\cmidrule(lr){4-5}\cmidrule(l){6-7}
Среднее поле & \(\alpha\) & \(\alpha_2\) & \(\alpha\) & \(\alpha_2\)
 & \(\alpha\) & \(\alpha_2\)\\
\midrule
HA  & 8,926 & 19,421 & 8,902 & 19,424 & \(-0{,}023\) & \(+0{,}003\)\\
QHA & 8,789 & 21,308 & 8,754 & 21,319 & \(-0{,}035\) & \(+0{,}011\)\\
\bottomrule
\end{tabular}
\end{center}
В приближении среднего поля (MFA) переход к классической статистике
расширяет двухфазный интервал QHA при \(p=0\) на \(0{,}046\) п.п.
HA рассчитана при фиксированном \(V_0/N=17{,}28833\)~\AA$^3$/атом.
